%%%%%%%%%%%%%%%%%%%%%%%%%%%%%%%%%%%%%%%%%
% Long Lined Cover Letter
% LaTeX Template
% Version 2.0 (September 17, 2021)
%
% This template originates from:
% https://www.LaTeXTemplates.com
%
% Authors:
% Vel (vel@LaTeXTemplates.com)
%
% License:
% CC BY-NC-SA 4.0 (https://creativecommons.org/licenses/by-nc-sa/4.0/)
%
%%%%%%%%%%%%%%%%%%%%%%%%%%%%%%%%%%%%%%%%%

%----------------------------------------------------------------------------------------
%	PACKAGES AND OTHER DOCUMENT CONFIGURATIONS
%----------------------------------------------------------------------------------------

\documentclass[10pt]{article}

\usepackage{charter} % Use the Charter font

\usepackage[
	a4paper, % Paper size
	top=.5in, % Top margin
	bottom=.6in, % Bottom margin
	left=.8in, % Left margin
	right=.8in, % Right margin
	%showframe % Uncomment to show frames around the margins for debugging purposes
]{geometry}

\setlength{\parindent}{0pt} % Paragraph indentation
\setlength{\parskip}{1em} % Vertical space between paragraphs

\usepackage{graphicx} % Required for including images
\usepackage{natbib}
\usepackage[colorlinks=true, linkcolor=magenta, citecolor=magenta, urlcolor=magenta]{hyperref}
\usepackage{fancyhdr} % Required for customizing headers and footers

\fancypagestyle{firstpage}{%
	\fancyhf{} % Clear default headers/footers
	\renewcommand{\headrulewidth}{0pt} % No header rule
	\renewcommand{\footrulewidth}{1pt} % Footer rule thickness
}

\fancypagestyle{subsequentpages}{%
	\fancyhf{} % Clear default headers/footers
	\renewcommand{\headrulewidth}{1pt} % Header rule thickness
	\renewcommand{\footrulewidth}{1pt} % Footer rule thickness
}

\AtBeginDocument{\thispagestyle{firstpage}} % Use the first page headers/footers style on the first page
\pagestyle{subsequentpages} % Use the subsequent pages headers/footers style on subsequent pages

%----------------------------------------------------------------------------------------

\begin{document}

%----------------------------------------------------------------------------------------
%	FIRST PAGE HEADER
%----------------------------------------------------------------------------------------

%\includegraphics[width=0.35\textwidth]{logo.pdf} % Logo

\vspace{-1em} % Pull the rule closer to the logo

\rule{\linewidth}{1pt} % Horizontal rule

\bigskip\bigskip % Vertical whitespace

%----------------------------------------------------------------------------------------
%	YOUR NAME AND CONTACT INFORMATION
%----------------------------------------------------------------------------------------

\hfill
\begin{tabular}{l @{}}
	\today \bigskip\\ % Date
\href{https://mamintoosi.github.io/}{Mahmood Amintoosi}\\
Division of Computer Science,\\
Department of Applied Mathematics \&\\
\href{https://gta-lab.github.io/}{Graph Theory and its Applications Lab} \\
Faculty of Mathematical Sciences,\\ Ferdowsi University of Mashhad,
IRAN
\end{tabular}

\bigskip % Vertical whitespace

Dear Editor,

\bigskip % Vertical whitespace

%----------------------------------------------------------------------------------------
%	LETTER CONTENT
%----------------------------------------------------------------------------------------

Thank you for the opportunity to revise our manuscript, ``Graph Structure Learning for Traffic Prediction'', for the \textit{International Journal of Data Science and Analytics}. We are grateful to the reviewers for their constructive comments. We have carefully addressed every point in the revised manuscript and in the accompanying point-by-point response letter.

Traditional graph-based traffic forecasting methods rely on a predefined adjacency built from physical road proximity, which need not match the dependency structure in the observed speeds. This work estimates the adjacency from the traffic data themselves using DAGMA and evaluates whether a learned graph improves GCN and T-GCN forecasting on SZ-Taxi and Los-loop.

The main changes in this revision are:
\begin{itemize}
    \item \textbf{Clearer constructions.} Contemporaneous and multi-lag DAGMA fits are specified separately. The earlier temporal reading of the contemporaneous DAG has been removed; learned graphs are interpreted as statistical associations, not causal maps.
    \item \textbf{Stronger controls.} We added matched edge-budget sparsity controls, an edge-budget sweep, a capacity check, structure figures, and an identity reference so that gains cannot be attributed only to sparsity or to removing a dense physical graph.
    \item \textbf{Statistical reporting.} Main results are five-seed means with sample standard deviations and paired per-seed consistency checks.
    \item \textbf{Scope and limitations.} Dataset dependence (including a 15-minute Los-loop control), offline DAGMA cost, and backbone scope are stated explicitly. Abstract percentages and causal wording were corrected.
\end{itemize}

On Los-loop, lag-specific learned graphs used in a lag-aligned way in T-GCN reduce RMSE by up to $43.0\%$ relative to the physical adjacency across horizons 1--4; matched-budget controls show that this gain is not explained by sparsity alone. On SZ-Taxi the improvements over the physical graph are much smaller.

The work remains aligned with the journal's scope in traffic prediction and graph-based knowledge discovery. A preliminary version was presented at a national conference \citep{Amintoosi1403aimc55-GSL}; the present manuscript is a substantial extension covering both GCN and T-GCN, contemporaneous and multi-lag constructions, multi-seed evaluation, and dedicated controls. Code and experiment scripts are available at \url{https://github.com/mamintoosi/TGCN-GSL-PyTorch}.

We believe the revised paper is clearer, better controlled, and more accurately stated. We look forward to your editorial decision.

\bigskip % Vertical whitespace

Sincerely,

\bigskip

Mahmood Amintoosi

\bibliographystyle{elsarticle-harv}
\bibliography{../MyReferences}

\end{document}
